%
% 1-bernoulli.tex
%
% (c) 2026 Prof Dr Andreas Müller
%
\section{Integration rationaler Funktionen nach Bernoulli
\label{buch:ratint:section:bernoulli}}
\kopfrechts{Integration rationaler Funktionen nach Bernoulli}%
Die Aufgabe, eine Stammfunktion zu finden, gehört zu den zentralen 
Themen jeder Analy\-sis-Grundvorlesung.
Eine besondere Rolle spielt darin der Fall einer rationalen
Funktion als Integrand.
Für diesen Fall gibt es eine vollständige Lösung, mindestens
im Prinzip.
Dieser Abschnitt zeigt diese Lösung von Bernoulli und hebt vor
allem hervor, warum sie für ein Computeralgebrasystem nicht
geeignet ist.

%
% 11-problem.tex -- Problemstellung
%
% (c) 2026 Prof Dr Andreas Müller
%
\subsection{Problemstellung}
Die von Bernoulli gelöste Aufgabe ist, die Stammfunktion einer
rationalen Funktion zu finden, für die ausserdem angenommen wird,
dass eine Faktorisierung des Nenners in Linearfaktoren bekannt ist.
Der Fundamentalsatz der Algebra garantiert, dass eine solche
Faktorisierung immer exisiert, wenngleich er keine Methode
bereitstellt, die Faktorisierung zu finden.

\begin{aufgabe}
Gegeben ist eine rationale Funktion $f(z) \in \mathbb{C}(z)$ mit
komplexen Koeffizienten.
\index{rationale Funktion}%
Finde eine Stammfunktion $F$.
\index{Stammfunktion}
\end{aufgabe}

Die rationale Funktion wird als gekürzter Bruch
\[
f(z)
=
\frac{p(z)}{q(z)}
\qquad\text{und}\qquad
p(z),q(z)\in \mathbb{Q}(z)
\]
geschrieben, wobei $p(z)$ und $q(z)$ teilerfremd sind.
Es wird angenommen, dass der Nenner in ein Produkt
\[
q(z)
=
(z-\alpha_1)\cdot\ldots\cdot (z-\alpha_n)
=
\prod_{k=1}^n (z-\alpha_k)
\]
faktorisiert werden kann.

Bernoullis Strategie ist, für die rationale Funktion erst
ein Partialbruchzerlegung zu finden, was im
Abschnitt~\ref{buch:ratint:bernoulli:subsection:partialbruch}
durchgeführt werden soll.
Er schreibt den Integranden als Summe
\begin{align*}
f(z)
=
s(z)
+
\sum_{k=1}^n\sum_{l=1}^{n_k}\frac{A_k^l}{(z-\alpha_k)^l},
\end{align*}
wobei $s(z)$ ein Polynom und $\alpha_k$ die Nullstellen
des Nenners von $f(z)$ sind. 
Um eine Stammfunktion zu finden, müssen Stammfunktionen der
einzelnen Summanden gefunden werden.
Der Polynomterm ist dabei das kleinste Problem, da die einzelnen
Summanden $z^n$ die Stammfunktion $z^{n+1}/(n+1)$ haben.
Etwas mehr Arbeit geben die Integrale der Brüche, die aber ebenfalls
wohlbekannt sind und in
Gleichung~\eqref{buch:ratint:bernoulli:partial:eqn:stammfunktionen}
angegeben werden.


\input{chapters/060-ratint/12-partial.tex}
%
% 13-reell.tex -- Rationale Funktionen mit reellen Koeffizienten
%
% (c) 2026 Prof Dr Andreas Müller
%
\subsection{Rationale Funktionen mit reellen Koeffizienten}
Der Fundamentalsatz der Algebra garantiert eine Faktorisierung 
eines Polynoms mit komplexen Koeffizienten in Linearfaktoren
der Form $z-\alpha$, wobei $\alpha$ eine Nullstelle des Polynoms
ist.
\index{Fundamentalsatz der Algebra}%
Dies ist jedoch nur möglich, wenn komplexe Nullstellen zugelassen
werden.
\index{Nullstelle}%
Es ist jedoch einfach, Polynome anzugeben, die über den reellen Zahlen
nicht in Linearfaktoren zerlegt werden können.
Das Polynom $p(x)=x^2+1$ hat keine reellen Nullstellen.
Über den reellen Zahlen ist $p(x)$ ein irreduzibles Polynom.
Über den komplexen Zahlen ist jedoch $p(x)=(x+i)(x-i)$ eine
Faktorisierung in Linearfaktoren.
\index{Faktorisierung}%

%
% Faktorisierung in lineare und quadratische Faktoren
%
\subsubsection{Faktorisierung in lineare und quadratische Faktoren}
Für ein reelles Polynom 
$p(z) = a_0 + a_1z + \dots + a_{n-1}z^{n-1} + a_nz^n$
kann man also eine Faktorisierung in Linearfaktoren
nicht erwarten.
Betrachtet man das Polynom jedoch als ein Polynom mit komplexen
Koeffizienten, dann garantiert der Fundamentalsatz der Algebra eine
Faktorisierung
\[
p(z)
=
a_n
(z-\alpha_1) 
(z-\alpha_2) 
\cdot
\ldots
\cdot
(z-\alpha_n),
\]
wobei die $\alpha_k$ Nullstellen des Polynoms sind.
Es gilt also $p(\alpha_k)=0$ für alle $k$.

Das Polynom $p(x)$ hat reelle Koeffizienten, die $\bar{a}_k=a_k$
für alle $k$ erfüllen.
Daher kann man die Bedingung $p(\alpha_k)=0$ komplex konjugieren
und bekommt
\begin{align*}
0
=
\overline{p(\alpha_k)}
&=
\bar{a}_0
+
\bar{a}_1 \bar{\alpha}_k
+
\bar{a}_2 \bar{\alpha}_k^2
+
\ldots
+
\bar{a}_{n-1} \bar{\alpha}_k^{n-1}
+
\bar{a}_{n} \bar{\alpha}_k^{n}
\\
&=
a_0
+
a_1 \bar{\alpha}_k
+
a_2 \bar{\alpha}_k^2
+
\ldots
+
a_{n-1} \bar{\alpha}_k^{n-1}
+
a_{n} \bar{\alpha}_k^{n}
\\
&=
p(\bar{\alpha}_k).
\end{align*}
Es folgt, dass $\bar{\alpha}_k$ ebenfalls eine Nullstelle ist.
Falls $\alpha_k=\bar{\alpha}_k$ ist, ist $\alpha_k$ eine reelle
Nullstelle.
Falls $\alpha_k\ne \bar{\alpha}_k$ ist, dann ist $\bar{\alpha}_k$
eine der anderen Nullstellen.
Die Nullstellen eines reellen Polynoms sind als entweder reell,
oder wenn sie nicht reell sind, dann treten sie in Paaren von
konjugiert komplexen Nullstellen auf.

Ein Paar von konjugiert komplexen Nullstellen führt auf den Faktor
\begin{align}
(z-\alpha)(z-\bar{\alpha})
&=
z^2 -(\alpha+\bar{\alpha}) z + \alpha\bar{\alpha}.
\label{buch:ratint:bernoulli:reell:faktoren}
\intertext{Die Summe konjugiert komplexer Zahlen ist
\index{konjugiert}%
$(a+bi)+(a-bi) = 2a $, also das Doppelte des Realteils.
Daher ist der zu berechnende Faktor}
&=
z^2 -2\operatorname{Re}(\alpha)\, z + |\alpha|^2.
\notag
\end{align}
Sowohl $\operatorname{Re}\alpha$ und $|\alpha|^2$ sind reell.
Indem man die Paare konjugiert komplexer Zahlen zu quadratischen
Faktoren wie in \eqref{buch:ratint:bernoulli:reell:faktoren}
zusammenfasst, findet man eine Faktorisierung des Polynoms in
der Form
\begin{equation}
p(z)
=
a_n
(z-\alpha_1)
\cdot\ldots\cdot
(z-\alpha_m)
\cdot
(z^2 + \beta_1 z + \gamma_1)
\cdot \ldots \cdot
(z^2 + \beta_l z + \gamma_l),
\label{buch:ratint:bernoulli:eqn:reellefaktorisierung}
\end{equation}
wobei die $\alpha_1,\dots,\alpha_m$ die reellen Nullstellen sind
und die Faktoren $z^2 +\beta_k z +\gamma_k$ Paare von
konjugiert komplexen, aber nicht reellen Lösungen haben.
Damit die rechte Seite den richtigen Grad hat, muss $n=m+2l$ sein.

%
% Reelle Partialbruchzerlegung
%
\subsubsection{Reelle Partialbruchzerlegung}
In der Zerlegung
\eqref{buch:ratint:bernoulli:eqn:reellefaktorisierung}
können einzelne Faktoren auch mehrfach vorkommen.
Die Konstruktion der Partialbruchzerlegung produziert
daher für jeden Faktor Partialnenner in Potenzen nicht
grösser als in der Faktorisierung.
Dies ist im folgenden Satz zusammengefasst.

\begin{satz}
\label{buch:ratint:bernoulli:satz:reelle-partialbruchzerlegung}
Ist der Nenner $q(z)$ der rationalen Funktion $f(z)=p(z)/q(z)$ 
ein Polynom mit einer Faktorisierung der Form
\begin{equation}
q(z)
=
a_n
(z-\alpha_1)^{s_1}
\cdots
(z-\alpha_m)^{s_m}
(z^2+\beta_1 z + \gamma_1)^{r_1}
\cdots
(z^2+\beta_l z + \gamma_l)^{r_l}
\end{equation}
dann gibt es eine eindeutig bestimmte Partialbruchzerlegung der Form
\begin{equation}
\renewcommand{\arraycolsep}{2pt}
%\renewcommand{\arraystretch}{1.7}
\begin{array}{rcclcccl}
f(z)
&=&&
\displaystyle
\frac{A_1^1    }{ z-\alpha_1       }
&+& \ldots &+&
\displaystyle
\frac{A_1^{r_1}}{(z-\alpha_1)^{s_1}}
\\
&&&\hspace*{11pt}\vdots&&&&\hspace*{17pt}\vdots
\\[2pt]
&&+&
\displaystyle
\frac{A_m^{1}  }{ z-\alpha_m       }
&+& \ldots &+&
\displaystyle
\frac{A_m^{s_m}}{(z-\alpha_m)^{s_m}}
\\[10pt]
&&+&
\displaystyle
\frac{B_1^1    z + C_1^1    }{ z+\beta_1 z + \gamma_1       }
&+& \ldots &+&
\displaystyle
\frac{B_1^{r_1}z + C_1^{r_1}}{(z+\beta_1 z + \gamma_1)^{r_1}}
\\
&&&\hspace*{24pt}\vdots&&&&\hspace*{29pt}\vdots
\\[2pt]
&&+&
\displaystyle
\frac{B_l^1    z + C_1^l    }{ z+\beta_l z + \gamma_l       }
&+& \ldots &+&
\displaystyle
\frac{B_l^{r_l}z + C_l^{r_1}}{(z+\beta_l z + \gamma_l)^{r_l}}.
\end{array}
\label{buch:ratint:bernoulli:eqn:reelle-partialbruchzerlegung}
\end{equation}
Darin sind die Zahlen $A_i^k$, $B_i^k$ und $C_i^k$ reelle Zahlen.
\end{satz}

\begin{proof}
Die Form der Partialbruchzerlegung ist aus der früheren Diskussion,
die zu Satz~\ref{buch:analytisch:rational:satz:partialbruchzerlegung}
geführt hat, klar.
Es muss höchstens noch gezeigt werden, wie diese Zerlegung effizient
ermittelt werden kann.

Offenbar sind $s_i$ Zahlen $A_i^k$ zu bestimmen, insgesamt also
$s_1+\dots+s_m$ Unbekannte für Terme zu den Linearfaktoren.
Zu den quadratischen Faktoren sind je $s_i$ Unbekannte
\index{quadratischer Faktor}%
$B_i^k$ und $C_i^k$ zu bestimmen, also insgesamt
$2r_1+\dots+2r_l$ Unbekannte.
Die Gesamtzahl der Unbekannten ist $s_1+\dots+s_m+2r_1+\dots+2r_l=n$.


Multipliziert man die Gleichung
\eqref{buch:ratint:bernoulli:eqn:reelle-partialbruchzerlegung}
mit dem Nenner $q(z)$, dann bleibt auf der linken Seite das Polynom
$p(z)$ und auf der rechten ein Polynom, dessen Koeffizienten
Linearkombinationen der Unbekannten $A_i^k$, $B_i^k$ und $C_i^k$
sind.
Das Zählerpolynom hat Grad $<n$, es enthält also höchstens $n$
Koeffizienten.
Koeffizientenvergleich liefert dann insgesamt $n$ lineare Gleichungen
für die Unbekannten.
Das entstehende lineare Gleichungssystem hat also $n$ Gleichungen für
\index{lineares Gleichungssystem}%
$n$ Unbekannte.
\end{proof}

Die reelle Partialbruchzerlegung erlaubt die Berechnung einer
Stammfunktion durch Integration der einzelnen Terme.
Für die linearen Terme ist bereits die Stammfunktion
\[
\int\frac{dz}{z-\alpha}
=
\log (z-\alpha)
\]
bekannt.
Die Stammfunktion für die Terme zu den quadratischen Faktoren
des Nenners wird im nächsten Abschnitt bestimmt.

%
% Stammfunktionen für quadratische Faktoren
%
\subsubsection{Stammfunktionen für quadratische Faktoren}
Möchte man das Integral nur im Reellen berechnen, muss man auch 
Partialbrüche der Form
\begin{equation}
f(z)
=
\frac{Az+B}{(z^2 +\beta z + \gamma)^m}
\qquad\text{mit $m>1$}
\label{buch:ratint:bernoulli:eqn:quadrpotnenner}
\end{equation}
mit quadratischen Nennern integrieren.
Dies ist mit wohlbekannten Methoden der klassischen Analysis 
möglich, wir geben hier nur die Resultate an, die man zum Beispiel
in \cite[section 2.1]{buch:bronstein} dargestellt finden kann.

Für die Nennerpotenz $m=1$ findet man die Formel
\begin{align*}
\int 
\frac{Bz+C}{z^2+\beta z+\gamma}
\,dz
&=
\frac{B}{2} \log( z^2+\beta z+\gamma )
+
\frac{2C-\beta B}{\!\sqrt{4\gamma-\beta^2}}
\arctan\frac{2z+\beta}{\!\sqrt{4\gamma-\beta^2}}.
\end{align*}
Hier treten logarithmische Terme und ein Arkustangensterm auf.
Die Arkustangensfunktion kann ebenfalls durch Logarithmusfunktionen
in der Form $\arctan(z) =-\frac{i}2\log(z+i) +\frac{i}2\log(z-i)$
ausgedrückt werden.
Im vorliegenden Fall ist dies besonders naheliegend, weil die
beiden Nullstellen des Nenner konjugiert komplex sind.

Höhere Nennerpotenzen $m>1$ können mit der folgenden Rekursionsformel
reduziert werden:
\begin{align}
\int 
\frac{Bz+C}{(z^2+\beta z+\gamma)^m}
\,dz
&=
\frac{
(2C-\beta B)z+\beta C-2\gamma B
}{
(m-1)(4\gamma-\beta^2)(z^2+bz+c)^{m-1}
}
\notag
\\
&\qquad
+
\int \frac{
(2m-3)(2C-\beta B)
}{
(m-1)(4\gamma-\beta^2)(z^2+\beta z+\gamma)^{m-1}
}
\,dz.
\label{buch:ratint:reell:eqn:rekursion}
\end{align}
Die Rekursionsformel basiert im wesentlichen auf partieller Integration.
Sie wird sich im Abschnitt~\ref{buch:ratint:section:hermite}
als ein Spezialfall der Methode von Hermite herausstellen.

%
% Die allgemeine Form der Stammfunktion
%
\subsubsection{Die allgemeine Form der Stammfunktion}
Die vorangegangenen Berechnungen zeigen, dass das Integral einer
rationalen Funktion $f(z)$ sich immer als eine Summe von Termen
der folgenden Form darstellen lassen.
\begin{enumerate}
\item
Für jeden Linearfaktor $z-\alpha$ sind rationale Terme der Form
\[
\frac{A}{(z-\alpha)^m}
\]
möglich, wobei $m$ nicht grösser sein kann als die Potenz, mit der
der Faktor im Nenner von $f(z)$ vorkommt.

Ausserdem ist ein Term der Form
\[
A \log (z-\alpha)
\]
möglich.
\item
Für jeden quadratischen Faktor der Form $q(z)=z^2+\beta z+\gamma$
sind rationale Terme der Form $t(z) / q(z)^m$ möglich, wobei $t(z)$
ein lineares Polynom ist und $m$ nicht grösser sein kann als die
Potenz, in der $q(z)$ im Nenner von $f(z)$ vorkommt.

Ausserdem sind Terme der Form
\[
\log q(z)
=
\log (z^2 +\beta z + \gamma)
\]
und
\[
\arctan \zeta(z)
=
-\frac{i}{2}\log (\zeta+i)
+\frac{i}{2}\log (\zeta-i)
\qquad\text{mit }
\zeta(z)=\frac{2z+\beta}{\!\sqrt{4\gamma-\beta^2}}
\]
möglich.
Die Stammfunktion lässt sich also auch in diesem Fall durch
rationale Funktionen und eine Linearkombination von drei Logarithmen
ausdrücken.
Dies ist in Übereinstimmung mit den Erkenntnissen des Liouville-Prinzips,
welches in Abschnitt~\ref{buch:intalg:section:liouville} präsentiert wird.
\end{enumerate}

Die Methode von Bernoulli postuliert, dass die Argumente der
Logarithmusfunktionen die Linearfaktoren der Nullstellen des
quadratischen Nenners sind.
Es ist nicht offensichtlich, dass die Argument $\zeta+i$ und $\zeta-i$
tatsächlich solche Linearfaktoren des Polynoms $q(z)$ sind.
Wir multiplizieren die Faktoren daher aus und erhalten
\begin{align*}
(\zeta+i)(\zeta-i)
&=
\zeta^2+1
\\
&=
\frac{4z^2+4z\beta+\beta^2}{4\gamma-\beta^2}
+
1
=
\frac{4z^2+4z\beta+\beta^2+4\gamma-\beta^2}{4\gamma-\beta^2}
=
\frac{4z^2+4z\beta+4\gamma}{4\gamma-\beta^2}
\\
&=
\frac{4}{4\gamma-\beta^2}
(z^2+\beta z+\gamma).
\end{align*}
Die Faktoren $\zeta+i$ und $\zeta-i$ haben daher die gleichen
Nullstellen wie das Polynom $q(z)=z^2+\beta z+\gamma$.
Tatsächlich ist
\begin{align*}
0&=\zeta\pm i
=
\frac{2z+\beta\pm i\!\sqrt{4\gamma-\beta^2}}{\!\sqrt{4\gamma-\beta^2}}
=
\frac{2z+\beta\pm \!\sqrt{\beta^2 - 4\gamma}}{\!\sqrt{4\gamma-\beta^2}}
&&\Rightarrow&
z
&=
-\frac{\beta}2\mp \!\sqrt{\frac{\beta^2}{4}-\gamma}.
\end{align*}
Dies ist die Lösungsformel für die quadratische Gleichung
$z^2+\beta z+\gamma=0$.






%
% 14-prinzip.tex -- Die der Methode zugrundeliegenden Prinzipien
%
% (c) 2026 Prof Dr Andreas Müller
%
\subsection{Die der Methode zugrundeliegenden Prinzipien
\label{buch:ratint:bernoulli:subsection:prinzip}}
Es wurde bereits früher festgehalten, dass die dargestellte Methode
für ein Computeralgebrasystem nicht durchführbar ist, weil eine
Methode zur Bestimmung der Faktorisierung des Nenners in Linearfaktoren
oder quadratische reelle Faktoren fehlt.
\index{Computeralgebrasystem}%
Es muss daher eine alternative Vorgehensweise gefunden werden,
die mit einer leichter zu findenden Faktorisierung des Nenners
arbeiten kann.
Die Faktorisierung muss immer noch einige Eigenschaften erfüllen,
damit die einzelnen Schritte verallgemeinert werden können.
Die folgenden Schritt müssen auf die neue Art der Faktorisierung
verallgemeinert werden.

\begin{enumerate}
\item
Die Partialbruchzerlegung wurde dazu verwendet, die 
rationale Funktion in Term mit speziell einfachen Nennern
aufzuspalten.
\index{Partialbruchzerlegung}%
In jedem Term kamen höchstens Potenzen eines einzigen Faktors
des Nenners vor.
Wie in der Darstellung der Partialbruchzerlegung angedeutet
wurde, müssen die Faktoren teilerfremd sein.
\item
Für den Reduktionsschritt
\[
\frac{A}{(z-\alpha)^l}
\quad\to\quad
\frac{\tilde{A}}{(z-\alpha)^{l-1}}
\qquad\text{bzw.}\qquad
\frac{a(z)}{q(z)^l}
\quad\to\quad
\frac{\tilde{a}(z)}{q(z)^{l-1}}
\]
ist eine Rekursionsformel
nötig, die den Exponenten im Nenner reduziert.
\index{Reduktion}%
Wie sich später zeigen wird, beruht sie darauf, dass der
Zähler in der Form
\begin{equation}
p(z) = s(z) \cdot q(z) + t(z)\cdot q'(z)
\label{buch:ratint:printzip:eqn:euklid}
\end{equation}
geschrieben werden kann.
Das aus dem zweiten Summanden resultierende Integral kann
dann durch partielle Integration
\begin{align*}
\int \frac{p(z)}{q(z)^m}\,dz
&=
\int \frac{s(z)}{q(z)^{m-1}}\,dz
+
\int t(z)\cdot \frac{q'(z)}{q(z)^m} \,dz
\\
&=
\int \frac{s(z)}{q(z)^{m-1}}\,dz
+
\int t(z) \frac{d}{dz} \frac{1-m}{q(z)^{m-1}}\,dz
\\
&=
\int \frac{s(z)}{q(z)^{m-1}}\,dz
+
t(z)\frac{1-m}{q(z)^{m-1}}
-
\int \frac{(1-m)t'(z)}{q(z)^{m-1}}\,dz
\end{align*}
vereinfacht werden, was genau ein Rekursionsformel der verlangten
Art ist.
Die Zerlegung \eqref{buch:ratint:printzip:eqn:euklid}
ist genau dann möglich, wenn $q(z)$ und $q'(z)$ teilerfremd
sind, also keine gemeinsame Nullstelle haben.
\index{teilerfremd}%
\item
Es wird ein Algorithmus benötigt, der die Polynom $s(z)$ und
$t(z)$ auf effiziente Art und Weise finden kann.
Um dies im Allgemeinen zu schaffen, müssen wir in
Abschnitt~\ref{buch:ratint:section:euklid} den euklidischen Algorithmus
genauer studieren.% und ihn in Abschnitt
%Abschnitt~\ref{buch:ratint:section:hermite}
%auf quadratfreie Nenner $q(z)$ anwenden.
\item
Es müssen Integrale der Form 
\[
\int \frac{p(z)}{q(z)}\,dz
\]
gefunden werden, wobei $p(z)$ ein Polynom vom Grad $\deg p(z)<\deg q(z)$ ist.
Die Methode von Rothstein-Trager von
Abschnitt~\ref{buch:ratint:section:rothsteintrager}
wird dies möglich machen.
\end{enumerate}

Eine geeignete Faktorisierung, die immer gefunden werden kann,
ist die sogenannte quadratfreie Faktorisierung, die in Abschnitt 
\ref{buch:ratint:section:hermite}
zusammen mit der Reduktion der Integrale von rationalen
Funktionen mit Potenzen des Nennerfaktors vorgestellt wird.



